\subsection{Basics}
\label{section_basics}

The following is provided for Euclidean geometric algebra of two-dimensional space
(\emph{EGA2D}) using an orthormal basis \bv{1}, \bv{2}:
\begin{subequations}
    \begin{align}
        (\bv{1})^2 = \bv{1} \wedgedot \bv1 & = \bv{1} \cdot \bv1 = \bv{1}^2 = +1,
        \;\text{and}\; \bv{2}^2 = +1 \nonumber\\
        s_{2d} & = s \, \bm{1}
        \label{eq:scalar2d} \\
        v_{2d} & = v_1\bv{1} + v_2\bv{2} 
        \label{eq:vec2d} \\
        ps_{2d} & = ps \, \bv{12}  = ps \, \ps
        \label{eq:pscalar2d}
    \end{align}
\end{subequations}
equation (\ref{eq:scalar2d}) is the scalar part $s_{2d}$, equation (\ref{eq:vec2d})
contains the vector part $v_{2d}$, and equation (\ref{eq:pscalar2d}) contains the
pseudoscalar part $ps_{2d}$. The index $2d$ is typically omitted when clear from context.
The four basis elements are $\{\bm{1}, \bv{1}, \bv{2}, \bv{12}\}$. Using this basis a
multivector $M$ of $2d$-space is defined as
\begin{equation}
     M = s \, \bm{1} + v_1\bv{1} + v_2\bv{2} + ps \, \bv{12}
    \label{eq:mvec2d}   
\end{equation}
where equation (\ref{eq:mvec2d}) contains three parts: the scalar part $s \, \bm{1}$
(basis element is the scalar \textbf{1}; if \textbf{1} is not shown, it is implicitly
assumed for scalar values), the vector part $v$ (basis elements \bv{1} and \bv{2}) and the
pseudoscalar part $ps \, \bv{12} = ps \, \ps$ (basis element is $I_{2d} =
\bv{12}$, which is sometimes written as $\ps$ to show its character as
pseudoscalar of this space. $I_{2d}$ is a bivector in two-dimensional Euclidean space). \\

For Euclidean geometric algebra of three-dimensional space (\emph{EGA3D}) using an
orthonormal basis \bv{1}, \bv{2}, \bv{3} there are:
\begin{subequations}
    \begin{align}
        (\bv{1})^2 = \bv{1} \wedgedot \bv1 & = \bv{1} \cdot \bv1 = \bv{1}^2 = +1,
        \bv{2}^2 = +1, \;\text{and}\; \bv{3}^2 = +1 \nonumber\\ 
        s_{3d} & = s \, \bm{1}
        \label{eq:scalar3d} \\
        v_{3d} & = v_1\bv{1} + v_2\bv{2} + v_3\bv{3} 
        \label{eq:vec3d} \\ 
        B_{3d} & = b_{23}\bv{23} + b_{31}\bv{31} + b_{12}\bv{12} 
        \label{eq:bivec3d} \\ 
        ps_{3d} & = ps \, \bv{123}  = ps \, \ps
        \label{eq:pscalar3d}
    \end{align}
\end{subequations}
equation (\ref{eq:scalar3d}) is the scalar part $s_{3d}$, equation (\ref{eq:vec3d}) is the
vector part $v_{3d}$, equation (\ref{eq:bivec3d}) is the bivector part $B_{3d}$, and
equation (\ref{eq:pscalar3d}) is the pseudoscalar part $ps_{3d}$. The index $3d$ is
typically omitted when clear from context. Compared to the $2d$-case it is obvious, that
all parts depend on context, specifically on the dimensionality of the modelled space, and
thus need to be defined and treated accordingly.\footnote{\emph{hint}: Since \Cpp is a
statically typed language, those types need to be well-defined and distinguishable from
each other. This is especially important for the different scalar and pseudoscalar types.
The ga library uses type tagging extensively for that purpose.} The basis elements are
$\{\bm{1}, \bv{1}, \bv{2}, \bv{3}, \bv{23}, \bv{31}, \bv{12}, \bv{123}\}$. Using this
basis a multivector $M$ of $3d$-space is defined as
\begin{equation}
    M = s \, \bm{1} + v_1\bv{1} + v_2\bv{2} + v_3\bv{3} 
    + b_{23}\bv{23} + b_{31}\bv{31} + b_{12}\bv{12} + ps \, \bv{123}
    \label{eq:mvec3d}  
\end{equation}
where equation (\ref{eq:mvec3d}) contains four parts: the scalar component $s \, \bm{1}$
(basis element is the scalar \textbf{1}), the vector part $v$ (basis elements \bv{1},
\bv{2} and \bv{3}), the bivector part $B$ (basis elements \bv{23}, \bv{31} and \bv{12})
and the pseudoscalar part $ps \, \bv{123} = ps \, \ps$ (basis element is $I_{3d}
= \bv{123}$, which is sometimes written as $\ps$ to show its character as
pseudoscalar of this space. $I_{3d}$ is a trivector in three-dimensional Euclidean space).
\\

Projective geometric algebra of two-dimensional Euclidean space (\emph{pga2dp}) is
modelled by embedding it in a three-dimensional space. Using an orthonormal basis \bv{1},
\bv{2}, \bv{3} there are following basis vectors:
\begin{subequations}
    \begin{align}
        (\bv{1})^2 = \bv{1} \wedgedot \bv1 & = \bv{1} \cdot \bv1 = \bv{1}^2 = +1,
        \bv{2}^2 = +1, \;\text{and}\; \bv{3}^2 = 0  \nonumber\\ 
        s_{2dp} & = s \, \bm{1}
        \label{eq:scalar2dp} \\
        v_{2dp} & = v_1\bv{1} + v_2\bv{2} + v_3\bv{3} 
        \label{eq:vec2dp} \\ 
        B_{2dp} & = b_{23}\bv{23} + b_{31}\bv{31} + b_{12}\bv{12} 
        \label{eq:bivec2dp} \\ 
        ps_{2dp} & = ps \, \bv{321}  = ps \, \ps
        \label{eq:pscalar2dp}
    \end{align}
\end{subequations}
The most obvious difference is the basis vector~\bv{3} that squares to zero and the
pseudoscalar of that space that is chosen to be \bv{321}. The trivector \bv{321} squares
to $0$, because of it's component \bv{3} contained as a factor. Equation
(\ref{eq:scalar2dp}) is the scalar part $s_{2dp}$, equation (\ref{eq:vec2dp}) is the
vector part $v_{2dp}$, equation (\ref{eq:bivec2dp}) is the bivector part $B_{2dp}$, and
equation (\ref{eq:pscalar2dp}) is the pseudoscalar part $ps_{2dp}$. The index $2dp$ is
typically omitted when clear from context. The eight basis elements are $\{\bm{1}, \bv{1},
\bv{2}, \bv{3}, \bv{23}, \bv{31}, \bv{12}, \bv{321}\}$. Using this basis a multivector $M$
of two-dimensional projective space ($2dp$) is defined as
\begin{equation}
    M = s \, \bm{1} + v_1\bv{1} + v_2\bv{2} + v_3\bv{3} 
    + b_1\bv{23} + b_2\bv{31} + b_3\bv{12} + ps \, \bv{321}
    \label{eq:mvec2dp}  
\end{equation}
where equation (\ref{eq:mvec2dp}) contains four parts: the scalar component $s \, \bm{1}$
(basis element is the scalar \textbf{1}), the vector part $v$ (basis elements \bv{1},
\bv{2} and \bv{3}), the bivector part $B$ (basis elements \bv{23}, \bv{31} and \bv{12})
and the pseudoscalar part $ps \, \bv{321} = ps \, \ps$ (basis element is
$I_{2dp} = \bv{321}$, which is sometimes written as $\ps$ to show its character
as pseudoscalar of this space. It is a trivector in two-dimensional Euclidean projective
space that squares to zero). \\

Projective geometric algebra of three-dimensional Euclidean space (\emph{pga3dp}) is
modelled by embedding it in a four-dimensional space. Using an orthonormal basis \bv{1},
\bv{2}, \bv{3}, \bv{4} there are following basis vectors:
\begin{subequations}
    \begin{align}
        (\bv{1})^2 = \bv{1} \wedgedot \bv1 & = \bv{1} \cdot \bv1 = \bv{1}^2 = +1,
        \bv{2}^2 = +1, \bv{3}^2 = +1,\;\text{and}\; \bv{4}^2 = 0 \nonumber\\ 
        s_{3dp} & = s \, \bm{1}
        \label{eq:scalar3dp} \\
        v_{3dp} & = v_1\bv{1} + v_2\bv{2} + v_3\bv{3} + v_4\bv{4} 
        \label{eq:vec3dp} \\ 
        B_{3dp} & = b_{41}\bv{41} + b_{42}\bv{42} + b_{43}\bv{43} +
                    b_{23}\bv{23} + b_{31}\bv{31} + b_{12}\bv{12} 
        \label{eq:bivec3dp} \\ 
        t_{3dp} & = t_{423}\bv{423} + t_{431}\bv{431} + t_{412}\bv{412} + t_{321}\bv{321} 
        \label{eq:trivec3dp} \\ 
        ps_{3dp} & = ps \, \bv{1234}  = ps \, \ps
        \label{eq:pscalar3dp}
    \end{align}
\end{subequations}
Now we have a basis vector~\bv{4} that squares to zero. The bivector has six components,
and there is a trivector $t$. The pseudoscalar \bv{1234} is a quadvector that squares to
$0$, because of it's component \bv{4} contained as a factor. Equation (\ref{eq:scalar3dp})
is the scalar part $s_{3dp}$, equation (\ref{eq:vec3dp}) is the vector part $v_{3dp}$,
equation (\ref{eq:bivec3dp}) is the bivector part $B_{3dp}$, equation (\ref{eq:trivec3dp})
is the trivector part $t_{3dp}$, and equation (\ref{eq:pscalar3dp}) is the pseudoscalar
part $ps_{3dp}$. The index $3dp$ is typically omitted when clear from context. The sixteen
basis elements of $pga3d$ are: $\{\bm{1}, \bv{1}, \bv{2}, \bv{3}, \bv{4}, \bv{41},
\bv{42}, \bv{43}, \bv{23}, \bv{31}, \bv{12}, \bv{423}, \bv{431}, \bv{412}, \bv{321},
\bv{1234}\}$. A multivector $M$ of three-dimensional projective space ($3dp$) is defined
as
\begin{equation}
    \begin{split}
    M & = s \, \bm{1} \\
      & + v_1\bv{1} + v_2\bv{2} + v_3\bv{3} + v_4\bv{4} \\
      & + b_{41}\bv{41} + b_{42}\bv{42} + b_{43}\bv{43} 
        + b_{23}\bv{23} + b_{31}\bv{31} + b_{12}\bv{12} \\
      & + t_{423}\bv{423} + t_{431}\bv{431} + t_{412}\bv{412} + t_{321}\bv{321} \\
      & + ps \, \bv{1234}
    \label{eq:mvec3dp}
    \end{split}
\end{equation}
where equation (\ref{eq:mvec3dp}) contains five parts: the scalar component $s \, \bm{1}$
(basis element is the scalar \textbf{1}), the vector part $v$ (basis elements \bv{1},
\bv{2}, \bv{3}, \bv{4}), the bivector part $B$ (basis elements \bv{41}, \bv{42}, \bv{43},
\bv{23}, \bv{31}, \bv{12}), the trivector part $t$ (basis elements \bv{423}, \bv{431},
\bv{412}, \bv{321}) and the pseudoscalar part $ps \, \bv{1234} = ps \, \ps$
(basis element is $I_{3dp} = \bv{1234}$, which is sometimes written as $\ps$ to
show its character as pseudoscalar of this space. It is a quadvector in three-dimensional
Euclidean projective space that squares to zero). \\

It is important to note that it is seldomly required to use a full multivector in the
respective algebra. For most applications it is sufficient to use a blade or $k$-vector
exclusively, which represents one grade of the multivector, e.g. a vector or a bivector.
This reduces the effort in terms of memory and computational requirements significantly --
typically reduced to a level that is required to to solve the underlying application
problem anyway. \\ 

The complement operation used in the library is defined
in~\cite{Lengyel_pga-illuminated:2024}. It is uniquely determined with respect to the
outer product (not with respect to the geometric product, as commonly used for dualization
by many autors working with GA). This is essential to generate consistent signs of
expressions, operators, complements and duals. The left complement $\ubar{u}$ and the
right complement $\bar{u}$ are defined as
\begin{subeqnarray}
    % use \label for full label
    % use \slabel for sublabel
    \ubar{u} \wedge u & = & I_n \slabel{eq:lcmpl_def} \\
    u \wedge \bar{u}  & = & I_n \slabel{eq:rcmpl_def}
\end{subeqnarray}
with $I_n  = \bv{1} \wedge \bv{2} \wedge \ldots \wedge \bv{n}$ as the pseudoscalar of the
$n$-dimensional algebra, and the wedge symbol $(\wedge)$ as the operator symbol of the
outer product. The effect of the complement is to turn basis elements not contained in the
object $u$ into the constituting basis elements of the complement of the object $\ubar{u}$
or $\bar{u}$. Connecting the object and its complement with the wedge products creates an
object that fills the available space in the sense of requiring all basis elements of the
space to represent it. The full space itself is represented by the pseudoscalar of that
space, since the pseudoscalar $I_n$ is formed by the product of all $n$ basis vectors of
the space. An implication of the definition is that the complement ($\ubar{u}$ or
$\bar{u}$) is orthogonal to $u$. \\

It can be shown that the left and right complements fulfill
\begin{equation}
    \bar{\ubar{a}} = a
    \label{eq:double_cmpl}
\end{equation}
thus the left and right complement operations are inverses to each other. This is valid
independent of the sequence of application, independent of the grade of the argument $a$,
and independent of the dimension of the space modelled in the algebra. Based on this it is
possible to define so-called regressive products using rules analog to deMorgan's rules in
boolean algebra as follows (definitions can be found
in~\cite{Lengyel_pga-illuminated:2024} and in \cite{Browne_Grassmann-Algebra_Vol1:2012}):
\begin{subeqnarray}
    % use \label for full label
    % use \slabel for sublabel
    a \vee b & \equiv & \underline{\bar{a} \wedge \bar{b}} \slabel{eq:rwdg_std} \\
    a \vee b & \equiv & \overline{\ubar{a} \wedge \ubar{b}} \slabel{eq:rwdg_alt}
\end{subeqnarray}
The same complement is applied to both arguments first, the resulting expressions are
combined with the wedge product, and the inverse complement operation is used to transform
the outcome into the final result, the regressive product. The inverted wedge symbol
$(\vee)$ is used for the regressive wedge product. Both definitions~(\ref{eq:rwdg_std})
and~(\ref{eq:rwdg_alt}) are equivalent. Equation~(\ref{eq:rwdg_std}) is used to derive
expressions for regressive products in \path{ga/ga_prdxpr}. \\

There are some additional formulas for complements which are useful for working with GA
expressions and thus are provided here for reference. Taking the right and left
complements of~(\ref{eq:rwdg_std}) and~(\ref{eq:rwdg_alt}) we end up with
\begin{subeqnarray}
    % use \label for full label
    % use \slabel for sublabel
    \overline{a \vee b}  & = & \overline{\underline{\bar{a} \wedge \bar{b}}}
    \quad  =  \quad \bar{a} \wedge \bar{b}
    \slabel{eq:intermediate_std} \\
    \underline{a \vee b} & = & \underline{\overline{\ubar{a} \wedge \ubar{b}}}
    \quad  =  \quad  \ubar{a} \wedge \ubar{b}
    \slabel{eq:intermediate_alt}
\end{subeqnarray}
Substituting $a = \ubar{u}$ and $b = \ubar{v}$ into~(\ref{eq:intermediate_std}) and $a =
\bar{u}$ and $b = \bar{v}$ into~(\ref{eq:intermediate_alt}) and exchanging the left and
right hand sides of the equations we end up with
\begin{subeqnarray}
    % use \label for full label
    % use \slabel for sublabel
    u \wedge v & = & \overline{\ubar{u} \vee \ubar{v}} \slabel{eq:wdg_std} \\
    u \wedge v & = & \underline{\bar{u} \vee \bar{v}} \slabel{eq:wdg_alt}
\end{subeqnarray}
which shows that the outer and the regressive outer product can be expressed by each
other. \\

The complement operation is used to define a unique dualization operation that works for
cases with non-degenerate as well as for cases with degenerate metric, like in projective
geometric algebra $PGA$. This is the main advantage compared to the usual prodedure of
dualization by multiplying the object with the pseudoscalar (or its inverse) using the
geometric product, since the latter breaks down for spaces with degenerate metrics. \\

The metric is defined by the dot product of the basis vectors as
\begin{equation}
    \mathfrak{g}_{ij} \equiv \bv{i} \cdot \bv{j}
\end{equation}
It is a square matrix that is the identity matrix $I_n$ for $EGA$, since all basis vectors
square to $\bv{i} \cdot \bv{i} = +1$, and $\bv{i} \cdot \bv{j} = 0$ for $i \neq j$. It is
invertible in this case. However, for $PGA$ one basis vector squares to zero, and thus
there is one zero on the main diagonal of the corresponding matrix of the metric, which is
therefore not invertible. It is called degenerate for that reason. \\

The extended metric $\mathbf{G}$ is a so-called exomorphism matrix that fulfills
\begin{equation}
    \mathbf{G}(A \wedge B) = \mathbf{G}(A) \wedge \mathbf{G}(B)
    \label{eq:exomorphism}
\end{equation} for any multivector $A$ and $B$ (see~\cite{Lengyel_pga-illuminated:2024},
pp. 61ff.). \\
Regarding $\mathbf{G}(u) = \mathbf{G} u$ as a unary operation on $u$ we can define the
corresponding regressive operation by first taking the complement of the argument and then
applying the opposed complement on the result. This defines the regressive metric
$\mathbf{\mathbb{G}}$ (named anti-metric by Lengyel). It is defined as
\begin{equation}
    \mathbf{\mathbb{G}} u \equiv \underline{\mathbf{G} \overline{u}}
    = \overline{\mathbf{G} \underline{u}}
\end{equation}
$\mathbf{G}$ and $\mathbf{\mathbb{G}}$ are related by 
\begin{equation}
    \mathbf{G} \mathbf{\mathbb{G}} = det(\mathfrak{g}) I_n
\end{equation}
Thus, if all basis vectors square to $\pm 1$, we have $\mathbf{G} \mathbf{\mathbb{G}} =
\pm I_n$, and both are invertible. However, if one basis vector squares to zero, we get
$\mathbf{G} \mathbf{\mathbb{G}} = 0$, which means that none of them is invertible, i.e. we
have a degenerate metric.\\

For $\mathbf{\mathbb{G}}$ there is the relation
\begin{equation}
    \mathbf{\mathbb{G}}(A \vee B) = \mathbf{\mathbb{G}}(A) \vee \mathbf{\mathbb{G}}(B)
\end{equation}
that is complementary to equation (\ref{eq:exomorphism}). \\

In $PGA$ the product $\mathbf{G} u$ defines the bulk of of an object $u$. It is defined as
\begin{equation}
    \bulk{u} = \mathbf{G} u
\end{equation}
The bulk only contains components that \emph{do not} contain the projective dimension
(\bv{3} for $pga2dp$ or \bv{4} for $pga3dp$) as a factor. It contains positional
information on $u$ relative to the origin. Zero bulk means that $u$ contains the origin.\\

In $PGA$ the product $\mathbf{\mathbb{G}} u$ defines the weight of of an object $u$. It is
defined as
\begin{equation}
    \weight{u} = \mathbf{\mathbb{G}} u
\end{equation}
The weight selects components that \emph{do} contain the projective dimension (\bv{3} for
$pga2dp$ or \bv{4} for $pga3dp$) as a factor. It contains directional information on $u$,
i.e. its attitude and direction independent of its position. Zero weight means that $u$ is
contained in the horizon, i.e. it is infinitely far away.\\

This effective split for each object into bulk and weight in projective algebras can be
written as
\begin{equation}
    u = \bulk{u} + \weight{u}
\end{equation}
This split is also relevant for the corresponding norms
\begin{equation}
    \nrm{u} = \bulknrm{u} + \weightnrm{u} = \sqrt{A \cdot A} + \sqrt{A \circ A}
\end{equation}
where the left hand side is called the geometric norm and the right hand side consists of
the sums of the bulk and weight norms. The bulk norm returns a scalar value of the
corresponding algebra, whereas the weight norm provides a pseudoscalar value of the
respective algebra. Consequently, the geometric norm is an object consisting of two
distinct parts, defined here as a dual number in the corresponding algebra. When we apply
unitization to geometric objects we devide the object in such a way that the magnitude of
the weight norm becomes a value of one. Both, bulk and weight norms, have a positive
argument under the square root due to the definition of the inner product and the
regressive inner product.\\

Using the metric and the right complement the right bulk dual \rightbulkdual{A} is defined
as
\begin{equation}
    \rightbulkdual{A} = \overline{\mathbf{G}A} = \overline{\bulk{A}}
    \label{eq:rightbulkdual}
\end{equation}
where $A$ is an arbitrary multivector, $\mathbf{G}$ the extended metric and $\mathbf{G}A$
a matrix-vector-product between them. \rightbulkdual{A} is also called the Hodge dual. The
multiplication with the metric happens before the application of the complement
operation.\\

There is also a corresponding left dual $\leftbulkdual{A}$ operation using the metric and
the left complement. It is defined as
\begin{equation}
    \leftbulkdual{A} = \underline{\mathbf{G}A} = \underline{\bulk{A}}
    \label{eq:leftbulkdual}
\end{equation}
The Hodge star (symbolized by a filled star $\star$) is used as dualization operator in
both cases. The lower or upper position symbolizes the left and right application. The
difference between the hodge dual and the complement is in both cases that multiplying the
multivector with the extended metric happens \emph{before} taking the corresponding
complement. However, this does not lead to different results in spaces where the metric is
the identity matrix. It only yields different results in spaces with a more involved
metric, e.g. like in projective spaces with a non-invertible (degenerate) metric as we
have in $PGA$. \\

The effect of applying the metric before taking the complement is equivalent to applying
the complement operation on the bulk part of the object $u$ exclusively, and not on the
full object $u$. Of course, the dualization can be applied to the weight part as well.
This operation (symbolized by an unfilled star~``$\smallstar$'') uses the regressive
extended metric and is defined for the right weight dual \rightweightdual{A} using the
right complement as
\begin{equation}
    \rightweightdual{A} = \overline{\mathbf{\mathbb{G}}A} = \overline{\weight{A}}
        \label{eq:rightweightdual}
\end{equation}
for an arbitrary multivector $A$. The corresponding left weight dual \leftweightdual{A} is
defined as
\begin{equation}
    \leftweightdual{A} = \underline{\mathbf{\mathbb{G}}A} = \underline{\weight{A}}
        \label{eq:leftweightdual}
\end{equation}
using the left complement on the weight part of $A$.\\

Relations to the geometric product and the regressive geometric product are given by
\begin{subeqnarray}
    \rightbulkdual{A} & = &  \revtilde{A} \wedgedot I_n = \revtilde{A} \wedgedot
    \ps \slabel{eq:hodge_right} \\
    \leftbulkdual{A} & = & I_n \wedgedot \revtilde{A} = \ps \wedgedot
    \revtilde{A} \slabel{eq:hodge_left} \\
    \rightweightdual{A} & = &  \undertilde{A} \veedot \bm{1} \slabel{eq:weight_right} \\
    \leftweightdual{A} & = & \bm{1} \veedot \undertilde{A} \slabel{eq:weight_left}
\end{subeqnarray}
with $I_n = \bv{1} \wedge \bv{2} \wedge \ldots \wedge \bv{n} = \ps$ as the
pseudoscalar of the $n$-dimensional space. $\revtilde{A}$ is the reversion operation, and
$\undertilde{A}$ is the regressive reversion operation applied to $A$ (e.g. with $A =
\bv{1} \wedge \bv{2}$ the reverse is $\revtilde{A} = \bv{2} \wedge \bv{1}$), and
``$\wedgedot$'' and ``$\veedot$'' are the symbols for the geometric and regressive
geometric products. \\

Following identities link the outer product and the inner product. They hold for operands
of same grade: 
\begin{subeqnarray}
    A \wedge B^{\star} & = & (A \cdot B) I_n, \quad \text{when \path{gr(A)} =
    \path{gr(B)}} \\
    A_{\star} \wedge B & = & (A \cdot B) I_n, \quad \text{when \path{gr(A)} =
    \path{gr(B)}}
\end{subeqnarray}

The influence of the metric for the dualization operation can be seen in its effect for
the different algebras which also depends on the dimension of the modelled space:
\begin{itemize}
\item $EGA2D$: $\;\:A_{\star} = \text{\path{left_dual(A)}}$, $A^{\star} =
\text{\path{right_dual(A)}}$
\item $EGA3D$: $\;\:A_{\star} = \text{\path{dual(A)}}$, $A^{\star} =
\text{\path{dual(A)}}$
\item $PGA2DP$: $A_{\star} = \text{\path{bulk_dual(A)}}$, $A^{\star} =
\text{\path{bulk_dual(A)}}$ 
\item $PGA2DP$: $A_{\smallstar} = \text{\path{weight_dual(A)}}$, $A^{\smallstar} =
\text{\path{weight_dual(A)}}$
\item $PGA3DP$: $A_{\star} = \text{\path{left_bulk_dual(A)}}$, $A^{\star} =
\text{\path{right_bulk_dual(A)}}$
\item $PGA3DP$: $A_{\smallstar} = \text{\path{left_weight_dual(A)}}$, $A^{\smallstar} =
\text{\path{right_weight_dual(A)}}$
\end{itemize}
In even-dimensional spaces ($EGA2D$, $PGA3DP$) we have to distinguish left and right
operations, while in odd-dimensional spaces ($EGA3D$, $PGA2DP$) we don't, since they
produce the same result. The ``left'' operations are typically used on the left-hand side
of the binary operation and the ``right'' operations on the right-hand side of the product
operator. The only deviation from that pattern occurs for the definition of the regressive
operations, as defined in~(\ref{eq:rwdg_std}) or (\ref{eq:rwdg_alt}). \\

In a two-dimensional algebra, like $EGA2D$, the effect of the right dual of a vector $u$
is to turn the vector by $+90$°, i.e. into the orientation of $\bv{12}$, while the left
dual turns the vector by $-90$°, i.e. into the direction against the orientation of
$\bv{12}$. This statement is valid for both right- and left-handed coordinate systems. \\

The multiplicative identity of the geometric product is the scalar $\bm{1}$:
\begin{equation}
    \bm{1} \wedgedot a = a \wedgedot \bm{1} = a
    \label{eq:wedgedot_identity} 
\end{equation}
The geometric product combines the dimensions that are present in an object
representation, e.g. for $pga2dp$:
\begin{eqnarray}
    \label{eq:gpr_combine_what_is_there}
    \bv{1} \wedgedot \bv{1} = 1 \nonumber\\
    \bv{2} \wedgedot \bv{2} = 1 \\
    \bv{3} \wedgedot \bv{3} = 0 \nonumber
\end{eqnarray}

The multiplicative identity of the regressive geometric product is the pseudoscalar
$\ps = I_n$ of the respective algebra:
\begin{equation}
    \ps \veedot a = a \veedot \ps = a
    \label{eq:veedot_identity} 
\end{equation}
The regressive geometric product combines the dimensions that are \underline{not} present
in an object
representation, e.g. for $pga2dp$:
\begin{eqnarray}
    \label{eq:rgpr_combine_what_is_not_there}
    \overline{\bv{1}} \veedot \overline{\bv{1}} = \ps \nonumber\\
    \overline{\bv{2}} \veedot \overline{\bv{2}} = \ps \\
    \overline{\bv{3}} \veedot \overline{\bv{3}} = 0 \nonumber
\end{eqnarray}
where eqs.~(\ref{eq:rgpr_combine_what_is_not_there}) can be shown to be equivalent to
eqs.~(\ref{eq:gpr_combine_what_is_there}) by taking the left complement of both sides and
applying the complement definition~(\ref{eq:rwdg_std}) to it. \\

The inner product between multivectors $A$ and $B$ is defined as
\begin{equation}
    A \cdot B = A^{T} \mathbf{G} B
    \label{eg:inner_product_metric}
\end{equation}
with $\mathbf{G}$ as the extended metric and matrix multiplication on the right hand side.
It satisfies the identity
\begin{equation}
    A \cdot B  = \grprj{A \wedgedot \tilde{B}}{0} = \grprj{ B \wedgedot \tilde{A} }{0}
    \label{eq:inner_product_link_to_gpr}
\end{equation}
where $\wedgedot$ stands for the geometric product within the grade projection operator
\grprj{}{k} for grade $k$. \\
\\

The contraction is explicitly defined as
\begin{subequations}
    \begin{align}
    A \lcontr B & = A \ll B = A_{\star} \vee B
    \label{eq:lcontr} \\
    B \rcontr A & = B \gg A = B \vee A^{\star}
    \label{eq:rcontr}
    \end{align}
\end{subequations}
with $\star$ denoting the Hodge dual which is formed by the respective complement
operation (in spaces of even dimension the left or right complement respectively, and in
spaces of odd dimension the complement function regardless of the side of the operand).
For blades A and B, they satisfy
\begin{subequations}
    \begin{align}
    A \lcontr B &  = A \ll B = \grprj{B \wedgedot \tilde{A}}{gr(B)-gr(A)}
    \label{eq:lcontr_gpr} \\
    B \rcontr A & = B \gg A = \grprj{\tilde{A} \wedgedot B}{gr(B)-gr(A)}
    \label{eq:rcontr_gpr}
    \end{align}
\end{subequations}
and fulfill the requirement that $A \lcontr B = A \rcontr B = A \cdot B$ whenever $A$
and $B$ have the same grade. In this case the contractions reduce to the inner product. \\


The geometric product between and vector $v$ and a blade $B$ is defined in terms of
interior and exterior products, i.e. the right contraction $\rcontr$, the left contraction
$\lcontr$ and wedge product $\wedge$ as follows:
\begin{subequations}
    \begin{align}
    v \wedgedot B & =  B \rcontr v + v \wedge B
    \label{eq:gpr_Brc} \\
    B \wedgedot v & =  v \lcontr B + B \wedge v
    \label{eq:gpr_Blc}
    \end{align}
\end{subequations}
or sometimes expressed using the right shift ($\gg$) or left shift ($\ll$) operators for
the contractions (as implemented in the ga library)
\begin{subequations}
    \begin{align}
    v \wedgedot B & = B \gg v + v \wedge B
    \label{eq:gpr_Bshrc} \\
    B \wedgedot v & =  v \ll B + B \wedge v
    \label{eq:gpr_Bshlc}
    \end{align}
\end{subequations}
or sometimes by exclusively using the operations of the exterior algebra (wedge product
``$\wedge$'', bulk dual ``$\star$'' and regressive wedge product ``$\vee$'')
\begin{subequations}
    \begin{align}
    v \wedgedot B & = B \vee \rightbulkdual{v} + v \wedge B
    \label{eq:gpr_Bshrc_exterior} \\
    B \wedgedot v & = \leftbulkdual{v} \vee B + B \wedge v
    \label{eq:gpr_Bshlc_exterior}
    \end{align}
\end{subequations}
as an alternative. Each one is equivalent. \\

Every vector $v$ can be decomposed into parts $v_{\parall}$ parallel to the $k$-blade $B$
and $v_{\perp}$ perpendicular to $B$, such that $v = v_{\parall} + v_{\perp}$. For
equation {(\ref{eq:gpr_Bshrc})}: $B \gg v$ is a $(k-1)$-blade. If $v_{\parall} \neq 0$
then $B \gg v$ represents a subspace of $B$. $v \wedge B$ is a $(k+1)$-blade. If
$v_{\perp} \neq 0$ then $v \wedge B$ represents $span(v,B)$.

For arguments of equal grade the contractions reduce to the dot product, so that one can
write specifically for two vectors $a$ and $b$:
\begin{equation}
    a \wedgedot b =  a \cdot b + a \wedge b
\end{equation}

\cite{Lengyel_transwedge_product:2025} has shown that the geometric product can be derived
by using operations from Grassmann algebra exclusively. These are the wedge product, the
complement and the duality operations. This demonstrates that the geometric product is not
a fundamental operation in its own right, but can be derived from more fundamental basic
operations. \\

For that purpose Lengyel introduced the transwedge product of order $k$ by following
definition:
\begin{equation}
    a \underset{k}{\doublewedge} b = \sum_{c \, \in \, \mathcal{B}_k}
    (\ubar{c}\vee a) \wedge (b \vee \rightbulkdual{c})
    \label{eq:twdg_orig}
\end{equation}
where $\mathcal{B}_k$ is the set of all basis elements of order $k$. For $EGA3D$ it means
for example $\mathcal{B}_0 = \{1\}$, $\mathcal{B}_1 = \{\bv{1}, \bv{2}, \bv{3}\}$,
$\mathcal{B}_2 = \{\bv{23}, \bv{31}, \bv{12}\}$, and $\mathcal{B}_3 = \{\bv{123}\}$. From
these expressions the geometric product can be formed by
\begin{equation}
    a \wedgedot b = \sum_{k \, = \, 0}^{n} (-1)^{k(k-1)/2} a \underset{k}{\doublewedge} b
    \label{eq:grp_twdg}
\end{equation}

For $k=0$ the formula yields the outer product (wedge product). For higher values of $k$
other products are generated which sum up to the full geometric product. Thus the
transwedge products helps to decompose the geometric product. The most interesting for
applications are the wedge product resulting for $k=0$ and the transwedge product of order
$k=1$. \\

The regressive wedge product in equation~(\ref{eq:twdg_orig}) can be replaced by the wedge
product using the complement-based definitions given in equations~(\ref{eq:double_cmpl})
to (\ref{eq:wdg_alt}). That means for example that equation~(\ref{eq:twdg_orig}) can be
transformed to the equivalent expression
\begin{equation}
    a \underset{k}{\doublewedge} b = \sum_{c \, \in \, \mathcal{B}_k}
    (\underline{c \wedge \bar{a}}) \wedge
    (\underline{\bar{b} \wedge \overline{\rightbulkdual{c}}})
    \label{eq:twdg_wdg}
\end{equation}
expressing it by using the wedge product exclusively. \\

Since every product can by assigned a corresponding regressive product in analogy to
definition~(\ref{eq:rwdg_std}) or~(\ref{eq:rwdg_alt}), the transwedge
product~(\ref{eq:twdg_orig}) also has a regressive form:
\begin{equation}
    a \underset{k}{\doublevee} b \equiv
    \overline{\ubar{a} \underset{k}{\doublewedge} \ubar{b}} =
    \sum_{c \, \in \, \mathcal{B}_k} \overline{
    (\ubar{c}\vee \ubar{a}) \wedge (\ubar{b} \vee \rightbulkdual{c})}
    \label{eq:rtwdg_orig}
\end{equation}
which in turn can be expressed in terms of the wedge product exclusively as:
\begin{equation}
    a \underset{k}{\doublevee} b = \sum_{c \, \in \, \mathcal{B}_k}  \overline{
    (\underline{c\wedge a}) \wedge (\underline{b \wedge \overline{\rightbulkdual{c}}})}
    \label{eq:rtwdg_wdg}
\end{equation}
Equations~(\ref{eq:twdg_wdg}) and (\ref{eq:rtwdg_wdg}) are used for calculation of the
product tables of the (regressive) transwedge products of order $k$ and the resulting
(regressive) geometric products in \path{ga/ga_prdxpr}. The transwedge product is very
useful to show that the geometric product can be derived from primitives of Grassmann
algebra. When (regressive) transwedge product expressions are calculated, it can be shown
that each transwedge product and the resulting coefficient expressions are equivalent to
specific combinations of wedge-, commutator-, dot-products, and their regressive
counterparts. \\

\newpage